%! TEX root = ../main.tex

\clearpage

\section{Các yêu cầu giao tiếp bên ngoài}

\subsection{Giao diện người sử dụng}

\textit{<Mô tả các đặc điểm luận lý (logic) của giao diện giữa sản phẩm phần mềm và người sử dụng. 
Phần này có thể bao gồm các ảnh màn hình mẫu, các chuẩn giao diện người sử dụng đồ họa 
(GUI) hay các hướng dẫn mẫu mà chúng phải được tuân theo, các ràng buộc về cách bố trí màn 
hình, các chức năng (ví dụ: trợ giúp) và các nút chuẩn mà chúng xuất hiện trong mọi màn hình, 
các phím tắt, các chuẩn hiển thị thông báo lỗi, v.v. Xác định các thành phần của phần mềm mà 
chúng cần các giao diện người sử dụng.>}

\subsection{Giao tiếp phần cứng}

\textit{<Mô tả các đặc điểm luận lý và vật lý của từng giao tiếp giữa sản phẩm phần mềm và các thành 
phần phần cứng của hệ thống. Phần này có thể bao gồm các loại thiết bị được hỗ trợ, trạng thái 
của sự tương tác điều khiển và dữ liệu giữa phần mềm và phần cứng và các giao thức giao tiếp 
được sử dụng.>}

\subsection{Giao tiếp phần mềm}
\textit{<Mô tả các nối kết giữa sản phẩm phần mềm này và các thành phần phần mềm cụ thể khác (tên 
và phiên bản), bao gồm cơ sở dữ liệu, hệ điều hành, công cụ, thư viện, các thành phần thương 
mại được tích hợp. Định danh các thành phần (item) dữ liệu hay các thông điệp đi vào và đi ra 
khỏi hệ thống và mô tả mục đích của từng item. Mô tả các dịch vụ được cần đến và trạng thái của 
các giao tiếp. Xác định dữ liệu sẽ được chia sẻ giữa các thành phần phần mềm, v.v.>}

\subsection{Giao tiếp truyền thông tin}
\textit{<Mô tả các yêu cầu có liên quan tới bất cứ chức năng truyền thông tin nào được cần bởi sản 
phẩm này, bao gồm thư điện tử, trình duyệt web, các giao thức truyền thông tin của máy chủ của 
mạng, các dạng điện tử, v.v. Định nghĩa định dạng thông điệp. Định danh bất cứ chuẩn truyền 
thông tin nào mà chúng sẽ được sử dụng, chẳng hạn như FTP, HTTP. Xác định các vấn đề về mã 
hóa và bảo mật truyền thông, tỷ lệ truyền dữ liệu và các kỹ thuật đồng bộ.> }

\subsection{Các ràng buộc về thực thi và thiết kế}
\textit{<Mô tả các vấn đề mà nhà phát triển phải lưu ý. Chúng bao gồm: các chính sách hợp tác hay điều 
tiết; các giới hạn phần cứng (các yêu cầu về thời gian và bộ nhớ); các giao diện với những ứng 
dụng khác; các cơ sở dữ liệu, công cụ, kỹ thuật cụ thể sẽ được sử dụng; các yêu cầu ngôn ngữ; 
các giao thức giao tiếp; v.v>}
